\chapter{Arquitectura}
\label{chap:arquitectura}

\section{Enfoque arquitectonico implementado}
\label{sec:arquitectura-enfoque}

La plataforma usa una arquitectura Django organizada por aplicaciones de dominio. El proyecto central en \archivo{config/} concentra configuracion, URL raiz, ASGI y WSGI. Las aplicaciones locales viven en \archivo{apps/} y separan responsabilidades funcionales. Esta estructura permite ubicar rapidamente donde se atiende una solicitud, donde se valida informacion, donde se persiste el estado y donde se ejecuta la regla de negocio.

La arquitectura no implementa una API separada ni una aplicacion SPA. La interfaz se basa en renderizado server-side con Django Templates y JavaScript puntual para interacciones publicas e internas. El panel de torneo agrega AJAX sobre vistas Django tradicionales para guardar resultados, actualizar participantes y refrescar contenido. Esta combinacion reduce complejidad de infraestructura frontend y mantiene la regla de negocio en servicios Python, pero exige disciplina en templates, atributos \variable{data-*} y contratos de respuesta.

\begin{figure}[H]
  \centering
  \fbox{\parbox{0.88\textwidth}{\centering Marcador de diagrama: arquitectura general del sistema. Fuente prevista: DIA-01 en \archivo{planificacion/06_inventario_diagramas.md}.}}
  \caption{Marcador para diagrama de arquitectura general.}
  \label{fig:arquitectura-general}
\end{figure}

\section{Organizacion por aplicaciones}
\label{sec:arquitectura-apps}

\begin{longtable}{p{3.3cm}p{4cm}p{6.2cm}}
\toprule
Aplicacion & Responsabilidad & Impacto arquitectonico \\
\midrule
\archivo{apps.common} & Modelo abstracto \clase{TimeStampedModel} & Centraliza campos temporales reutilizables y evita repetir \variable{created_at} y \variable{updated_at} en modelos de dominio. \\
\archivo{apps.core} & Sitio publico, reglas y patrocinadores & Mantiene la entrada publica separada de registro, comunicaciones y panel interno. \\
\archivo{apps.participants} & Registros, participantes, asistencia, instituciones, equipos y miembros & Conecta el flujo publico con el estado operativo del torneo mediante sincronizacion hacia equipos. \\
\archivo{apps.tournament} & Ediciones, reglas, competencias, grupos, batallas, estados, historial, recomendaciones y panel & Contiene la logica critica de competencia y concentra la mayor sensibilidad de mantenimiento. \\
\archivo{apps.invitations} & Plantillas, destinatarios, lotes, logs, DOCX, PDF y correo & Aisla comunicaciones y reduce el riesgo de mezclar envio de correos con logica de torneo. \\
\bottomrule
\end{longtable}

La dependencia mas relevante ocurre entre participantes y torneo. \archivo{apps.participants} no termina su trabajo al guardar un registro; cuando se confirma asistencia, invoca servicios que crean o actualizan equipos y pueden inicializar competencia. \archivo{apps.tournament}, a su vez, depende de \clase{Team} y \clase{TeamMember} para construir grupos, batallas y estados por division. Esta relacion explica por que el manual trata el registro como parte del flujo competitivo y no como un formulario aislado.

\begin{figure}[H]
  \centering
  \fbox{\parbox{0.88\textwidth}{\centering Marcador de diagrama: dependencias entre aplicaciones. Fuente prevista: DIA-03 en \archivo{planificacion/06_inventario_diagramas.md}.}}
  \caption{Marcador para diagrama de dependencias entre aplicaciones Django.}
  \label{fig:arquitectura-dependencias-apps}
\end{figure}

\section{Aplicacion del patron MTV}
\label{sec:arquitectura-mtv}

El proyecto aplica el patron Model-Template-View de Django con una capa adicional de servicios para la regla de negocio. Los modelos representan el estado persistente: registros, participantes, equipos, ediciones, competencias, grupos, batallas, estados e historial. Las vistas reciben solicitudes HTTP, aplican permisos, preparan formularios, llaman servicios y devuelven templates, redirecciones, JSON o fragmentos HTML segun el flujo. Los templates construyen la interfaz publica e interna, y el JavaScript del panel agrega comportamiento dinamico sobre esas respuestas.

La capa de servicios es decisiva en este proyecto. Si toda la regla de torneo quedara en vistas, cada endpoint podria replicar logica de clasificacion, sincronizacion o correccion. En cambio, \archivo{apps/tournament/services.py} concentra operaciones como \funcion{initialize_competition}, \funcion{set_group_qualifiers}, \funcion{set_battle_winner}, \funcion{materialize_repechage_stage}, \funcion{generate_manual_phase_proposal}, \funcion{confirm_manual_phase_proposal}, \funcion{sync_competition} y \funcion{save_final_podium}. Esto permite que distintas vistas usen la misma regla de negocio. El costo es que ese archivo requiere pruebas y revision rigurosa porque cualquier cambio puede impactar varios flujos.

\begin{figure}[H]
  \centering
  \fbox{\parbox{0.88\textwidth}{\centering Marcador de diagrama: patron MTV aplicado al proyecto con capa de servicios. Fuente prevista: DIA-02 en \archivo{planificacion/06_inventario_diagramas.md}.}}
  \caption{Marcador para diagrama del patron MTV aplicado a la plataforma.}
  \label{fig:arquitectura-mtv-servicios}
\end{figure}

\section{Flujo de peticion HTTP}
\label{sec:arquitectura-http}

La URL raiz del proyecto define cinco entradas principales:

\begin{itemize}
  \item \ruta{/admin/}: administracion Django.
  \item \ruta{/}: sitio publico de \archivo{apps.core}.
  \item \ruta{/registro/}: registro y confirmacion de participantes.
  \item \ruta{/comunicaciones/}: panel interno de comunicaciones.
  \item \ruta{/torneo/}: vista publica e interna del torneo.
\end{itemize}

Las rutas publicas de \archivo{apps.core} y \archivo{apps.participants} no exigen autenticacion. Esto permite consultar informacion y registrar equipos sin crear cuentas. La confirmacion de asistencia usa token UUID como mecanismo de identificacion del registro. En contraste, las rutas de \archivo{apps.tournament} para \ruta{/torneo/control/} y las rutas de \archivo{apps.invitations} bajo \ruta{/comunicaciones/} requieren autenticacion y validacion de organizador, definido como usuario \variable{is_staff} o \variable{is_superuser}.

El flujo HTTP interno sigue este patron: la URL resuelve una vista; la vista valida permisos y parametros; si hay formulario o accion POST, llama servicios de dominio; el servicio actualiza modelos y registra efectos; la vista responde con template, redireccion, JSON o HTML parcial. En operaciones AJAX, \archivo{static/js/control.js} envia encabezados \variable{X-CSRFToken} y \variable{X-Requested-With}; la vista procesa la solicitud y devuelve una respuesta que el cliente usa para confirmar, reintentar o refrescar el panel.

\begin{figure}[H]
  \centering
  \fbox{\parbox{0.88\textwidth}{\centering Marcador de diagrama: flujo de peticion HTTP y AJAX. Fuente prevista: DIA-05 y DIA-13 en \archivo{planificacion/06_inventario_diagramas.md}.}}
  \caption{Marcador para diagrama de flujo HTTP y operaciones AJAX.}
  \label{fig:arquitectura-flujo-http}
\end{figure}

\section{Persistencia}
\label{sec:arquitectura-persistencia}

La persistencia se implementa con Django ORM. En desarrollo local, el proyecto usa SQLite cuando \variable{DJANGO_DATABASE=sqlite} o cuando \variable{DEBUG=True} sin variables PostgreSQL. La configuracion impide usar SQLite con \variable{DEBUG=False}; si no se usa SQLite, exige variables PostgreSQL. Esta regla protege despliegues no locales contra una configuracion accidentalmente debil o incompleta.

Los modelos persistidos se agrupan por dominio. \archivo{apps.participants} mantiene registros de colegio y universidad, participantes, instituciones, equipos y miembros. \archivo{apps.tournament} mantiene ediciones, reglas, fases, competencias por division, grupos, entradas de grupo, batallas, participantes de batalla, estados de equipos e historial. \archivo{apps.invitations} mantiene plantillas, destinatarios, lotes y logs. La persistencia no solo guarda datos; tambien permite reconstruir estado operativo, historial de decisiones y resultados de comunicaciones.

\section{Logica de torneo}
\label{sec:arquitectura-logica-torneo}

La logica de torneo combina perfiles automaticos y decisiones manuales. Los perfiles automaticos definen estructuras segun cantidad de equipos: modo pendiente con menos de cuatro, configuraciones con grupos y final four para rangos bajos, estructuras con cuartos o ronda de 32 para rangos superiores y repechajes cuando corresponde. Este comportamiento permite responder a la cantidad real de equipos confirmados, que puede variar respecto a inscripciones iniciales.

Las fases manuales existen para resolver escenarios que no pueden tratarse como avance automatico simple. El sistema genera propuestas, las guarda temporalmente en \variable{competition.configuration}, valida firmas y participantes, permite reasignaciones y materializa batallas solo cuando la propuesta se confirma. Esa estrategia protege contra duplicados, grupos vacios o clasificaciones inconsistentes. El manual debe desarrollar estos contratos en capitulos posteriores porque son parte central del mantenimiento del sistema.

\section{Frontend y AJAX}
\label{sec:arquitectura-frontend-ajax}

El frontend publico usa \archivo{static/js/app.js} para tema, accesibilidad del boton de tema, comportamiento del header y animaciones de revelado. El panel interno usa \archivo{static/js/control.js}, que contiene comportamiento de mayor impacto operativo: modales de participante, guardado de resultados, seleccion de ganadores o clasificados, decision de fase, asistente manual, guardado multiple y refresco parcial.

La arquitectura frontend evita introducir un framework cliente adicional. La ventaja es que el estado principal permanece en Django y la interfaz se renderiza desde templates conocidos por el backend. El costo tecnico es el acoplamiento entre JavaScript y atributos \variable{data-*}. Si un template cambia el nombre de un atributo usado por \archivo{control.js}, una accion interna puede dejar de funcionar aunque el HTML siga renderizando. Por eso, los capitulos de frontend y JavaScript deben documentar esos contratos con precision.

\section{Administracion}
\label{sec:arquitectura-administracion}

El proyecto usa Django Admin como herramienta administrativa complementaria. \archivo{apps.participants/admin.py} registra administradores para instituciones, registros y equipos, e incluye acciones que pueden enviar solicitudes de asistencia o sincronizar registros confirmados hacia equipos. \archivo{apps.tournament/admin.py} registra modelos de edicion, reglas, fases, competencias, grupos, batallas, estados e historial. \archivo{apps.invitations/admin.py} permite revisar plantillas, destinatarios, lotes y logs.

La administracion aporta capacidad de inspeccion y correccion, pero tambien introduce riesgo operativo. Algunas acciones administrativas pueden disparar correos o modificar datos que alimentan el torneo. Por esa razon, el manual no trata Django Admin como un detalle secundario: queda incluido dentro de la arquitectura porque forma parte de las superficies reales de operacion y mantenimiento.

\section{Riesgos arquitectonicos conocidos}
\label{sec:arquitectura-riesgos}

La auditoria validada identifica riesgos que afectan la arquitectura:

\begin{itemize}
  \item \textbf{Concentracion de logica}: \archivo{apps/tournament/services.py} concentra muchas responsabilidades. Esto facilita encontrar la regla de negocio, pero exige pruebas focalizadas y revisiones cuidadosas.
  \item \textbf{Configuracion dinamica}: \variable{DivisionCompetition.configuration} almacena informacion operativa en JSON. La flexibilidad permite soportar propuestas manuales y configuracion de fases, pero requiere documentar claves esperadas.
  \item \textbf{Dependencia de LibreOffice}: la conversion DOCX a PDF depende de un binario externo disponible por PATH o \variable{LIBREOFFICE_BINARY}.
  \item \textbf{Acoplamiento JS-template}: \archivo{control.js} depende de atributos \variable{data-*} en templates internos.
  \item \textbf{Cobertura de pruebas limitada}: existen pruebas para distribuciones manuales y comunicaciones, pero faltan pruebas end-to-end de flujos completos.
  \item \textbf{Produccion no definida}: dominio, servidor, reverse proxy, certificados, backups y almacenamiento de media permanecen como \MarcadorProduccionPendiente.
\end{itemize}

Estos riesgos no invalidan la arquitectura. Indican donde debe concentrarse la validacion cuando se redacten los capitulos de modelo de datos, servicios, flujo de competencia, pruebas, despliegue y mantenimiento.
